\chapter{Conclusion}
\label{ch:future-work}
\epigraph{``Nothing I did before the age of seventy was worth noticing.''}{Hokusai, One Hundred Views of Mount Fuji}

\lettrine{W}{e} set out to show that peephole rewriting can be verified without taking either the compiler or the solver on trust.
Toward this, we built a framework to describe MLIR compiler IRs and their peephole rewrites in Lean,
and we built verified decision procedures to automatically discharge proof obligations that arise in this process.
%They are evaluated against the unverified state of the art, and none of them matches its raw speed.
%Our floating-point bitblaster trails Bitwuzla by a geometric-mean factor of $\SmtlibRandSpeedupBitwuzlaOverOurs\times$,
%and the fixed-width bitvector engine it is built on is itself within roughly an order of magnitude of Bitwuzla's.
%What that price buys is an end-to-end correctness guarantee that reaches down to Lean's kernel.
Together, our solvers provide evidence that the finite-domain reasoning which underpins compilers
can be carried out inside a proof assistant, at scale.

Looking back, I can see many open problems that remain. I'll describe the three frontiers I find most exciting:
nonlinear arithmetic over parametric bitvectors, a fully mechanized Bitwuzla, and tensor bitblasting.

\section{Nonlinear Bitvector Arithmetic via the 2-adic Integers}

The parametric solvers of this thesis (\cref{ch:mono-width,ch:multi-width})
decide the \emph{linear-bitwise} fragment of the parametric bitvector theory.
I have an unpublished sketch showing that \emph{nonlinear} arithmetic in the
parametric bitvector theory is decidable, by a reduction to the first-order
theory of the 2-adic integers, and, conversely, that integrating \emph{any}
bitwise operation into such a theory immediately leads to
undecidability.\footnote{%
\url{https://pixel-druid.com/articles/non-linear-theory-of-2-adics-does-not-mix-with-bitwise-operations}}
This implies that the nonlinear arithmetic fragment of the parametric
bitvector theory is (maximally) decidable: to write incomplete solvers for
the full theory, one must first build a complete solver for the nonlinear
arithmetic fragment, which can then be mixed with bitwise operations via
theory combination.

Toward this, I propose implementing a cylindrical algebraic decomposition (CAD) solver for the 2-adics.
While the decidability of the full first-order theory of the 2-adics has been known in theory since the
1960s~\citep{ax1966diophantine}, that result was not algorithmic. A CAD-style
algorithm based on cell decomposition does exist~\citep{quantifier-elimination-padic},
but, to the best of my knowledge, has never been implementation. I therefore propose to implement
such a solver, leveraging existing libraries such as FLINT~\citep{hart2013flint} for polynomial manipulation and root finding.
This also sets the stage for a fully verified reimplementation in Lean, extending
the reach of the verified solvers of this thesis beyond the linear-bitwise fragment.

\section{Fully Mechanized Bitwuzla}

The key missing piece of the verified solver is support for the theory of arrays.
Arrays can be implemented in terms of a larger body of theory, called \emph{ackermannization},
which reduces array reasoning to uninterpreted functions and bitvectors.
Uninterpreted functions are captured entirely by their congruence rule,
$x = y \implies f(x) = f(y)$. So, one creates new atoms called $f_x$ and $f_y$, which stand in for $f(x)$ and $f(y)$, and adds the axiom
$x = y \implies f_x = f_y$. This is a standard technique for adding support for uninterpreted functions to a SAT based theory.
However, since there are quadratically many pairs of terms,
doing this eagerly is not feasible.
Instead, one can do this \emph{incrementally}, only adding the axiom when the solver finds a counterexample to it.

Adding \emph{incremental solving} to \bvdecide{} would therefore enable much better
variants of ackermannization, and unlock support for the theory of
arrays. Together with the floating-point support of \cref{ch:floating-point},
this would bring the verified solver to theory parity with
Bitwuzla~\citep{niemetz2023bitwuzla} ($\qfbv{} + \qffp{} + \qfuf{}$),
which suffices to model almost all behaviour of realistic production languages such as C and C++.
Industry formalization efforts would be a natural testing ground for
evaluating whether such an extended solver can handle real verification
workloads. Bitwuzla's ackermannization and multiplication-abstraction
machinery is a natural starting point for a mechanization.

\section{Tensor Bitblasting}

Tensor bitblasting should be possible, and I'm going to sketch out the algorithms that I believe can be used to implement this.
The overall idea is to reduce the problem to thinking of tensor operations as a composition of``reshaping'' operations,
and ``elementwise'' operations, which together can be combined to decide the theory of tensors.
Firstly, for elementwise operations, FpSan~\citep{openai2026fpsan} gives floating-point level equivalence, and guarantees equivalence upto an exponential ring,
which is what ML practitioners care about anyway. It can be implemented as a bitvector solver, and I have a prototypical implementation of it in Lean.
Secondly, TensorRight~\citep{arora2025tensorright} is an algorithm for shape level equivalence.
The key idea is similar to that of our parametric bitvector solver,
that shows that for every shape-based equivalence, one can compute a finite upper bound on the dimensions of the tensors that must be checked.
Then, checking upto this cutoff establishes equivalence for all unbounded ranks.
Note that this does not immediately give a SAT-friendly encoding.
To then use the ideas to get a bitblastable encoding, we propose using (3) Linear layouts~\citep{zhou2026linear}, which showed how by making the mild assumption that all tensor sizes are powers of 2,
basically all tensor based shape operations can be encoded efficiently as linear algebra over $Z/2Z$. This of course provides extremely convenient bitvector encodings.
The three ingredients are in place, and the missing piece is to combine them correctly, and filling the holes in each theory with support from the others.

I expect the masking trick of our multi-width solver (\cref{ch:multi-width})
can be adapted to tensor shapes, since a symbolic dimension plays much the same role as a symbolic width,
and FpSan can be used to reduce the floating-point reasoning to bitvectors.
The payoff would be a tensor bitblaster that is fast, because it is powered by SAT,
and that transfers from bounded to unbounded checks, because it is built on TensorRight's theory,
and is trustworthy, because it can be implemented via bitblasting in a proof assistant.

\section*{Closing}

Looking back, for me, it was worth wondering why bitblasting works in practice.
I read it as a flavour of the bitter lesson~\citep{sutton2019bitter},
the observation from AI that search and learning are the only two paradigms that keep rewarding the compute we feed them.
Automated reasoning has its own incarnation of search in SAT solving,
which now routinely settles instances with hundreds of millions of variables~\citep{biere2024cadical,kissatSATComp2024},
and bitblasting~\citep{tseitin1983complexity} is the craft of encoding richer theories into this engine.
Moreover, the SAT community has figured out how to produce certificates that can be rechecked by a small kernel,
which is exactly what foundational trust demands.

It is also worth wondering whether the effort of implementing these procedures in Lean~\citep{demoura2021lean}
and proving them correct is worth it. I believe it is, and that formal verification will only grow in importance.
These are interesting times (as of September 2026): large language models~\citep{brown2020language,touvron2023llama}
produce impressive volumes of code, but cannot be trusted to get the details right.
A machine-checked specification, with machine-checked automation to discharge it is the clear path forward.

Overall, a large body of an optimizing compiler rests on peephole rewriting,
and the correctness of each rewrite is a small but crucial piece of the overall correctness argument.
The procedures that prove these rewrites en masse need not be taken on trust.
They can be implemented, proved correct, and run inside a proof assistant,
on the problems practitioners actually have.
Foundational proof and push-button automation need not be at odds,
and the price of foundational trust need not be too steep to pay.
